\section{The Most Partisan Parameter Configurations}
\label{sec:most-partisan}

\subsection{Defining Most-Partisan Configurations}

To identify the most partisan parameter settings, we need to understand
which direction each parameter favors.
From U.2's single-parameter sweeps:

\begin{table}[h]
\centering
\caption{Direction of partisan effect for each parameter.
         ``Rep-favorable'' means higher setting produces fewer Democratic seats;
         ``Dem-favorable'' means higher setting produces more Democratic seats.
         $\hat{\Delta}D$ per unit is the estimated marginal effect on $\Dnat$
         per unit change in the parameter.}
\label{tab:param-direction}
\begin{tabular}{lcccl}
\toprule
Parameter & Symbol & Rep-favorable end & Dem-favorable end & Effect direction \\
\midrule
$\ufactor$ & $u_{\rm factor}$ & 5.0\% ($\Dnat = 210.8$) & 0.1\% ($\Dnat = 211.2$) & Higher = less compact = Dem \\
$\acounty$ & $\alpha_{\rm county}$ & 10.0 ($\Dnat = 210.9$) & 1.0 ($\Dnat = 211.1$) & Higher = more county-split resistance = Rep \\
$T$ & $T$ & 100 ($\Dnat = 211.1$) & 600+ ($\Dnat = 211.0$) & Lower = less optimized = Dem \\
$\aswing$ & $a_{\rm swing}$ & 1.20 ($\Dnat = 210.9$) & 1.05 ($\Dnat = 211.1$) & Higher = more area imbalance = Rep \\
$\wvra$ & $w_{\rm vra}$ & 0.6 ($\Dnat = 210.9$) & 0.2 ($\Dnat = 211.1$) & Higher VRA = Rep (packs minority communities) \\
\bottomrule
\end{tabular}
\end{table}

\subsection{Most Republican-Favorable Configuration}

Setting all parameters to their most Republican-favorable values:
\[
  \theta_{\rm Rep} = (\ufactor = 5\%, \; \acounty = 10, \; T = 1{,}000,
  \; \aswing = 1.20, \; \wvra = 0.6).
\]
Note: $T = 1{,}000$ at the Republican-favorable end because the single-parameter
sweep shows no partisan effect from $T$ at $T \geq 200$; we include it at the
extreme for completeness but expect no marginal contribution from this choice.

\subsection{Most Democratic-Favorable Configuration}

Setting all parameters to their most Democratic-favorable values:
\[
  \theta_{\rm Dem} = (\ufactor = 0.1\%, \; \acounty = 1, \; T = 100,
  \; \aswing = 1.05, \; \wvra = 0.2).
\]

\subsection{Expected Additive Bound}

If the five parameters were independent and additively partisan, the expected
difference between $\theta_{\rm Rep}$ and $\theta_{\rm Dem}$ would be:
\begin{align*}
  \Delta \Dnat^{\rm additive} &=
  |210.8 - 211.2| + |210.9 - 211.1| + |211.1 - 211.0| \\
  &\quad + |210.9 - 211.1| + |210.9 - 211.1| \\
  &= 0.4 + 0.2 + 0.1 + 0.2 + 0.2 = 1.1 \text{ seats}.
\end{align*}
The empirical joint effect will be smaller due to saturation and correlation.
